\chapter{Conclusão}
\label{cap:conclusao}

O projeto exigia duas tarefas simultâneas: separar o canal de interesse dos canais vizinhos e proteger a decimação por 4 contra o dobramento espectral. Como os sinais eram originados de um sistema SDR e baseados em uma constelação I/Q, o filtro FIR passa-baixas projetado pelo método da janela de Kaiser cumpriu ambas as funções com precisão. Com 93 coeficientes e um atraso de grupo constante de 46 amostras, a solução atendeu a todos os três critérios de aprovação com bastante folga.

Uma situação observada ao responder os questionários foi que a atenuação de $60 dB$ presente na tabela de dados não era suficiente. O critério de aprovação exige que o interferente de $450 kHz$ fique $60 dB$ abaixo do canal útil, mas ele entra no sistema $15,56 dB$ acima dele. Somando essa diferença de entrada ao nível de rejeição exigido, a atenuação real necessária no filtro passa a ser de $75,56 dB$. Desta forma, um filtro projetado com exatamente $60 dB$ seria reprovado nesse requisito.

Na prática, a escolha pelo filtro FIR se mostra totalmente acertada quando analisada em detalhes: embora um filtro IIR elíptico de ordem 6 conseguisse resolver a seletividade com bem menos coeficientes, ele causaria uma variação de 24 amostras no atraso de grupo dentro do canal, borrando a constelação I/Q que o bloco seguinte precisa ler, enquanto o FIR simétrico garante fase linear por construção, mantendo um atraso rigorosamente constante em 46 amostras. Além disso, a arquitetura polifásica demonstra que o alto custo computacional do FIR é facilmente superado, pois ao dividir a filtragem em quatro subfiltros que já operam na taxa reduzida, o número de multiplicações por amostra de saída cai de 372 para 96 sem qualquer aproximação — apresentando uma diferença irrisória de apenas $1{,}5 \times 10^{-15}$ por erro de arredondamento —, o que equipara o custo do FIR polifásico ao de um IIR equivalente, com a grande vantagem de não distorcer a fase e de ser incondicionalmente estável.

Por fim, a comparação dos espectros antes e depois da decimação deixou visível o fenômeno que motivou todo o projeto. Sem filtro, o interferente de $450$~kHz reaparece em $-30$~kHz com $15{,}6$~dB acima do canal útil. Com o filtro aplicado antes do decimador, ele desaparece. A ordem das operações, filtrar primeiro, decimar depois é o que torna a redução de taxa uma operação segura. Foi muito interessante esse projeto pois me ajudou a entender melhor como funciona o Processamento de Sinais, e foi um desafio finaliza-lo e compreender e integrar todo o sistema deste o contexto apresentado, no final foi uma ótima experiência de aprendizado.